\begin{lemma}[the toric support function detects central restriction]
\label{lem:analytic-support-specialization}
\Verified\ Let \(M\in\operatorname{Pic}(N_0)\), let
\(\pi:Y_\epsilon\to N_0\) be the cusp covering, and let
\(\nu:E_0\to W\) be the normalization.  If the Appell--Humbert form of
\(M|_{F_t}\), \(t\ne0\), is \(E_M\), then the three independent degrees
of \(\nu^*(M|_W)\) determine an integral symmetric form
\(B:\mathbb Z^2\times\mathbb Z^2\to\mathbb Z\) satisfying
\[
 E_M\bigl((\lambda,0),(0,m)\bigr)
 =-B(B_0\lambda,m)
 \qquad(\lambda,m\in\mathbb Z^2).
\tag{5.17-A21}
\]
Consequently
\[
 E_M=0
 \quad\Longrightarrow\quad
 \nu^*(M|_W)\simeq\mathcal O_{E_0}.
\tag{5.17-A22}
\]
\end{lemma}

\begin{proof}
Put \(\mathcal M=\pi^*M\).  The identity
\(\pi\circ\Psi_\lambda=\pi\) gives \(\mathcal M\) its genuine
\(\Psi\)-equivariant descent structure.  In the notation of
\cref{prop:local-cusp-picard-obstruction}, order the six rays in the star
of the origin as
\[
 (r_1,r_2,r_3,r_4,r_5,r_6)
 =(e_1,e_2,e_2-e_1,-e_1,-e_2,e_1-e_2),
 \qquad C_i=E_0\cap E_{r_i}.
\tag{5.17-A23}
\]
The quotient identifies opposite conductor branches.  Therefore, for
\(A=\mathcal M|_{E_0}=\nu^*(M|_W)\),
\[
 d_i=\deg_{C_i}A,
 \qquad
 (\deg_{C_1}A,\ldots,\deg_{C_6}A)
 =(d_1,d_2,d_3,d_1,d_2,d_3).
\tag{5.17-A24}
\]

Define the integer-valued quadratic function
\[
 a(x,y)
 =(d_2+d_3)\binom{x}{2}+d_3xy+(d_1+d_3)\binom{y}{2}.
\tag{5.17-A25}
\]
Its polarization is
\[
 B(v,w)=a(v+w)-a(v)-a(w),
 \qquad
 [B]=
 \begin{pmatrix}
  d_2+d_3&d_3\\ d_3&d_1+d_3
 \end{pmatrix}.
\tag{5.17-A26}
\]
The locally finite divisor
\[
 D_a=\sum_{v\in\mathbb Z^2}a(v)E_v
\tag{5.17-A27}
\]
is Cartier: on every maximal chart, the three height-one ray generators
form an integral basis, so the three assigned coefficients are the values
of one integral character.  In the standard toric convention the associated
support function has ray values \(-a(v)\); the divisor coefficients here are
the values \(a(v)\), consistent with
\(\operatorname{div}(\chi^m)=\sum_\rho\langle m,u_\rho\rangle D_\rho\)
\cite[\S3.3, pp.~60--63; \S3.4, pp.~63--72]{Fulton1993}%
\cite[Propositions~4.1.1--4.1.2, p.~171; Theorems~4.2.8,
p.~179 and~4.2.12, p.~182]{CoxLittleSchenck2011}.
This is the same chartwise Cartier argument as
in (5.17ab), and it applies to an arbitrary integer-valued \(a\).

Direct substitution gives
\[
 (a(r_1),\ldots,a(r_6))
 =(0,0,d_2,d_2+d_3,d_1+d_3,d_1).
\tag{5.17-A28}
\]
Using \(C_i^2=-1\) and \(C_iC_{i\pm1}=1\), the degree vector of
\(\mathcal O(D_a)|_{E_0}\) is exactly the vector in (5.17-A24).  More
generally, restrict the principal identity
\(\operatorname{div}(t)=\sum_wE_w\) to \(E_v\).  Since \(E_v\) meets
precisely \(E_{v+r_i}\), it gives
\[
 \mathcal O(D_a)|_{E_v}
 \simeq
 \mathcal O_{E_v}\!\left(
   \sum_{i=1}^6(a(v+r_i)-a(v))C_{v,i}
 \right).
\tag{5.17-A29}
\]
The coefficient difference is
\[
 a(v+r_i)-a(v)=a(r_i)+B(v,r_i).
\tag{5.17-A30}
\]
The second summand in (5.17-A30) is the divisor of the toric character
with exponent \(B(v,-)\).  Thus every restriction in (5.17-A29) has the
same six degrees as \(\mathcal O(D_a)|_{E_0}\).

Because \(B_0\bar\Lambda=\mathbb Z^2\), the deck action is transitive on
the components \(E_v\).  The genuine descent structure on \(\mathcal M\)
therefore transports (5.17-A24), with its ordered six boundary directions,
to every \(E_v\).  It follows that
\[
 \mathcal R:=\mathcal M\otimes\mathcal O(-D_a)
 \quad\text{satisfies}\quad
 \mathcal R|_{E_v}\simeq\mathcal O_{E_v}
 \quad(v\in\mathbb Z^2).
\tag{5.17-A31}
\]
Indeed, the six toric curves detect \(\operatorname{Pic}(\mathrm{dP}_6)\):
in the blow-up basis \(H,E_1,E_2,E_3\), vanishing of the degrees on the
three exceptional curves kills the three \(E_i\)-coefficients, and
vanishing on any strict transform \(H-E_i-E_j\) kills the remaining
\(H\)-coefficient.  Since
\(H^1(\mathrm{dP}_6,\mathcal O)=0\), equality of those degrees implies
isomorphism, not merely numerical equivalence.

For \(b_\lambda=B_0\lambda\), the torus-translation part of
\(\Psi_\lambda\) preserves each toric prime divisor, while its fan part
translates the index by \(b_\lambda\).  Hence
\[
 \Psi_\lambda^*D_a-D_a
 =\sum_v(a(v+b_\lambda)-a(v))E_v
 =\operatorname{div}(\chi^{m_\lambda}),
\tag{5.17-A32}
\]
where
\[
 m_\lambda=(\ell_\lambda,k_\lambda),
 \qquad
 \ell_\lambda=B(b_\lambda,-),
 \qquad
 k_\lambda=a(b_\lambda).
\tag{5.17-A33}
\]
Thus multiplication by \(\chi^{m_\lambda}\) gives an individual
isomorphism
\(\Psi_\lambda^*\mathcal O(D_a)\simeq\mathcal O(D_a)\).
No cocycle assertion for these individual isomorphisms is needed: we use
them only to separate the horizontal winding of the genuine descent factor
of \(\mathcal M\).

For a ray \(\mathbb R_{\ge0}(v,1)\), its toric open set in
\(Y_\epsilon\) is
\[
 U_v\simeq\Delta_\epsilon\times(\mathbb C^*)^2,
 \qquad E_v^\circ=\{0\}\times(\mathbb C^*)^2.
\tag{5.17-A34}
\]
Indeed \((e_1,0),(e_2,0),(v,1)\) is an integral basis and \(t\) is the
character dual to the last vector; compare
\cite[\S2.1, pp.~28--29]{Fulton1993} and
\cite[Example~1.2.21, p.~31]{CoxLittleSchenck2011}.
The set \(U_v\) is Stein and retracts
onto \(E_v^\circ\).  By (5.17-A31),
\(c_1(\mathcal R|_{U_v})=0\); the exponential sequence and Cartan's
theorem~B therefore trivialize \(\mathcal R|_{U_v}\).  Choose a global
nowhere-zero frame \(s_v\) of \(\mathcal R|_{E_v}\).  A frame on \(U_v\) can be
multiplied by a unit pulled back from \((\mathbb C^*)^2\) so that its
restriction to \(E_v^\circ\) agrees with \(s_v|_{E_v^\circ}\); denote the
resulting frame by \(e_v\).

We also need to compare these frames on the generic torus.  The direction
of the toric affine-open inclusions is essential here.  If \(v,w\) are
adjacent vertices, choose a triangle \(\sigma=\{v,w,u\}\) containing them,
write \(z_v,z_w,z_u\) for the coordinates dual to its three height-one ray
generators, and put
\[
 S_\sigma:=U_\sigma\cap Y_\epsilon
 =\{(z_v,z_w,z_u)\in\mathbb C^3:
       |z_vz_wz_u|<\epsilon\}.
\tag{5.17-A34a}
\]
Because \(\rho_v<\sigma\), localization of the maximal affine chart gives
\begin{equation}
 U_{\rho_v}=D(z_wz_u)\subset U_\sigma,\qquad
 U_{\rho_w}=D(z_vz_u)\subset U_\sigma,\qquad
 U_{\rho_v}\cap U_{\rho_w}=D(z_vz_wz_u)=U_0.
\label{eq:analytic-toric-open-directions}
\end{equation}
Thus \(U_\sigma\) is not contained in \(U_{\rho_v}\cap U_{\rho_w}\), and
the coordinate discs in \(S_\sigma\) cannot be applied directly to the
ratio \(e_w/e_v\).

The space \(S_\sigma\) is nevertheless the required comparison chart.  It
is the closed graph
\[
 \{(z_v,z_w,z_u,t)\in\mathbb C^3\times\Delta_\epsilon:
   t=z_vz_wz_u\}
\]
in the Stein space \(\mathbb C^3\times\Delta_\epsilon\), and hence is
Stein.  Radial scaling of the three \(z\)-coordinates contracts it to the
origin.  The exponential sequence and Cartan's theorem~B therefore give
\(\operatorname{Pic}(S_\sigma)=0\).  Choose a frame \(e_\sigma\) of
\(\mathcal R|_{S_\sigma}\).

Let \(s_v\) be the global frame of
\(\mathcal R|_{E_v}\simeq\mathcal O_{E_v}\) used to normalize \(e_v\).
On \(U_{\rho_v}\cap Y_\epsilon\) put
\[
 A_{\sigma,v}=e_v/e_\sigma.
\]
At \(t=0\), this is the restriction to \((\mathbb C^*)^2\) of the
holomorphic unit
\[
 G_{\sigma,v}
 =\frac{s_v}{e_\sigma|_{E_v\cap S_\sigma}}
 \in\mathcal O^*(E_v\cap S_\sigma),
 \qquad E_v\cap S_\sigma=\{z_v=0\}\simeq\mathbb C^2.
\]
Consequently its winding on every horizontal cocharacter is zero at
\(t=0\).  Under
\(U_{\rho_v}\cap Y_\epsilon\simeq
\Delta_\epsilon\times(\mathbb C^*)^2\), that winding is an integer locally
constant in \(t\), and is therefore zero on every nearby nonzero fibre.
The same argument applies to \(A_{\sigma,w}\).  On their actual common
domain \(U_0\cap Y_\epsilon\),
\[
 e_w/e_v=A_{\sigma,w}/A_{\sigma,v},
\]
so this ratio has zero horizontal winding.  The one-skeleton of the
\(A_2\)-triangulation is connected, and multiplication along an edge path
now gives
\[
 \operatorname{wind}_m(e_w/e_v)=0
 \qquad(v,w,m\in\mathbb Z^2)
\tag{5.17-A34b}
\]
on every nonzero fibre.

Let
\(\theta^{\mathcal M}_\lambda:\Psi_\lambda^*\mathcal M\to\mathcal M\)
be the genuine descent isomorphism, let
\(\mu^+_\lambda:\Psi_\lambda^*\mathcal O(D_a)\to\mathcal O(D_a)\)
be multiplication by \(\chi^{m_\lambda}\) from (5.17-A32), and define the
inverse-dual isomorphism
\[
 \mu^-_\lambda:=((\mu^+_\lambda)^{-1})^\vee:
 \Psi_\lambda^*\mathcal O(-D_a)\longrightarrow\mathcal O(-D_a),
 \qquad s\longmapsto\chi^{-m_\lambda}s.
\]
The ordinary dual of \(\mu^+_\lambda\) has the opposite arrow; the
inverse-dual above is the map required by the pullback tensor product.  Then
\[
 \theta^{\mathcal R}_\lambda
 :=\theta^{\mathcal M}_\lambda\otimes\mu^-_\lambda:
 \Psi_\lambda^*\mathcal R\longrightarrow\mathcal R.
\]
Since \(\Psi_\lambda(U_v)=U_{v+b_\lambda}\), define
\(r_{\lambda,v}\in\mathcal O^*(U_v)\) by the typed identity
\begin{equation}
 \theta^{\mathcal R}_\lambda
   (\Psi_\lambda^*e_{v+b_\lambda})
 =r_{\lambda,v}e_v.
\label{eq:analytic-residual-descent-factor}
\end{equation}
At \(t=0\) it extends to an isomorphism between two trivial line bundles
on the compact surface \(E_v\).  Its factor in the chosen global frames is
therefore a nonzero constant.  The winding of
\(r_{\lambda,v}(t,-)\) around every horizontal one-parameter subgroup is
an integer locally constant in \(t\); at \(t=0\) it is zero, hence it is
zero for every nearby \(t\ne0\).

To express this isomorphism in one frame on the generic torus, put
\(q_{\lambda,v}=e_{v+b_\lambda}/e_v\).  Equation (5.17-A34b) gives zero
horizontal winding for \(q_{\lambda,v}\); because
\(\Psi_\lambda|_{Y_t}\) is multiplication by a fixed torus point, it
induces the identity on horizontal \(H_1\), and
\(\Psi_\lambda^*q_{\lambda,v}\) also has zero winding.  Relative to the
single frame \(e_v\), the residual multiplier is therefore
\[
 h_{\lambda,v}
 =\frac{r_{\lambda,v}}{\Psi_\lambda^*q_{\lambda,v}},
 \qquad
 \operatorname{wind}_m(h_{\lambda,v})=0.
\tag{5.17-A34c}
\]

On \(U_0\) let \(\mathbf1_a\) be the canonical toric frame of
\(\mathcal O(D_a)\), and put \(F_v=e_v\otimes\mathbf1_a\).  Equations
\cref{eq:analytic-residual-descent-factor} and (5.17-A32) give
\[
 \theta^{\mathcal M}_\lambda(\Psi_\lambda^*F_v)
 =h_{\lambda,v}\chi^{m_\lambda}F_v.
\]
On a nonzero fibre, \(\chi^{m_\lambda}=x^{\ell_\lambda}t^{k_\lambda}\).
With the section-descent convention (5.17am), the factor of automorphy is
therefore
\[
 \alpha_\lambda=(h_{\lambda,v}\chi^{m_\lambda})^{-1}.
\]
The residual factor has zero horizontal winding by (5.17-A34c), and hence
\begin{equation}
 \operatorname{wind}_m(\alpha_\lambda)
 =-\ell_\lambda(m)=-B(B_0\lambda,m).
\label{eq:analytic-factor-horizontal-winding}
\end{equation}

This winding is exactly the mixed Chern-class entry.  Pull \(F_v\) back to
the universal cover
\[
 p_t:\mathbb C^2\longrightarrow Y_t,
 \qquad \zeta\longmapsto(e^{2\pi i\zeta},t).
\]
The period indexed by \((\lambda,n)\in\mathbb Z^2\oplus\mathbb Z^2\)
acts by \(\zeta\mapsto\zeta+Z(s)\lambda+n\).  The horizontal factor is one,
because \(F_v\) is already defined on \(Y_t\).  Choose a logarithm
\(g_\lambda\in\mathcal O(\mathbb C^2)\) with
\[
 p_t^*\alpha_\lambda=e^{2\pi i g_\lambda}.
\]
For every horizontal period \(n\),
\[
 g_\lambda(\zeta+n)-g_\lambda(\zeta)
 =\operatorname{wind}_n(\alpha_\lambda).
\]
In the exponential connecting cocycle
\[
 c(l,l')=g_l(\zeta+l')+g_{l'}(\zeta)-g_{l+l'}(\zeta),
\]
take the horizontal logarithm to be zero and the logarithm for the combined
period \((\lambda,n)\) to be \(g_\lambda\).  Then
\[
 c((\lambda,0),(0,n))
 =\operatorname{wind}_n(\alpha_\lambda),
 \qquad
 c((0,n),(\lambda,0))=0.
\]
Therefore the alternating exponential connecting integer between
the vertical deck
translation indexed by \(\lambda\) and the horizontal loop indexed by
\(m\) is
\[
 E_M\bigl((\lambda,0),(0,m)\bigr)
 =-\ell_\lambda(m)
 =-B(B_0\lambda,m),
\]
which proves (5.17-A21).  This sign also recovers the explicit check in
\cref{prop:local-cusp-picard-obstruction}: for
\(a(v)=v^{\mathsf t}Hv\), one has \(B=2H\), so (5.17-A21) is the
cross-block \(-2B_0^{\mathsf t}H=-2P\), namely \(-2q\).

If \(E_M=0\), equation (5.17-A21) and the unimodularity of \(B_0\) force
\(B=0\).  The matrix in (5.17-A26) then gives
\(d_3=0\), \(d_2=0\), and \(d_1=0\).  Thus all six degrees of \(A\)
vanish, and the Picard detection argument following (5.17-A31) gives
\(A\simeq\mathcal O_{E_0}\).  This is (5.17-A22).
\end{proof}

